\begin{tikzpicture}[font=\small,>=stealth]
\draw[rounded corners=3pt,gray!45] (0,0) rectangle (14.2,12.2);

\node at (7.1,11.55) {先拿掉双曲线：只比较一条直线上的三个点};
\node at (7.1,10.25) {上图：中间点 $P$，两端在相反方向};
\draw[gray] (1.2,8.6)--(12.9,8.6);
\draw[blue!75!black,very thick] (6.2,8.6)--(10.2,8.6);
\draw[orange!85!black,very thick] (6.2,8.6)--(3.2,8.6);
\fill (10.2,8.6) circle(2pt) node[above=9pt] {$A$};
\fill (6.2,8.6) circle(2pt) node[above=9pt] {$P$};
\fill (3.2,8.6) circle(2pt) node[above=9pt] {$B$};
\draw[->,blue!75!black] (6.5,7.9)--(9.7,7.9);
\draw[->,orange!85!black] (5.9,7.9)--(3.7,7.9);
\node at (7.1,6.85) {$AB=PA+PB$，两段的水平位移方向相反。};
\node at (7.1,5.35) {下图：外部点 $P$，两端在同一方向};
\draw[gray] (1.2,3.8)--(12.9,3.8);
\draw[orange!85!black,very thick] (11.2,3.8)--(2.2,3.8);
\draw[blue!75!black,very thick] (11.2,3.8)--(7.2,3.8);
\fill (11.2,3.8) circle(2pt) node[above=9pt] {$P$};
\fill (7.2,3.8) circle(2pt) node[above=9pt] {$A$};
\fill (2.2,3.8) circle(2pt) node[above=9pt] {$B$};
\draw[->,blue!75!black] (10.8,3.1)--(7.7,3.1);
\draw[->,orange!85!black] (10.8,2.6)--(2.7,2.6);
\node at (7.1,1.6) {$AB=PB-PA$，两段的水平位移方向相同。};
\node at (7.1,.55) {此图只认位置和方向；倒数的加减还需要后面的焦半径关系。};
\end{tikzpicture}
